%%%%%%%%%%%%%%%%%%%%%%%%%%%%%%%%%%%%%%%%%%%%%%%%%%%%%%%%%%%%%%%%%%%%%%
% Lớp 11 | Hàm số lượng giác và phương trình lượng giác | Các công thức lượng giác
% Dạng 1D1?3-6: Nhận dạng tam giác
% Số câu: 2  (sinh tự động bởi tach_cau.py; từ nay thầy sửa trực tiếp file này)
%%%%%%%%%%%%%%%%%%%%%%%%%%%%%%%%%%%%%%%%%%%%%%%%%%%%%%%%%%%%%%%%%%%%%%

%%%%%%%%%%%%%%%% TRẢ LỜI NGẮN %%%%%%%%%%%%%%%%

% nguồn: Ok_TraLoiNgan_Lop11_Lan2
\begin{ex}%[1D1V3-6]%[ma:f0ad0358]
	Trong tam giác $ABC$ ta có $\cos A+\cos (B+C)+\tan \dfrac{A+B}{2}=k \cot \dfrac{C}{2}$. Khi đó $k$  bằng
	\shortans{1} 
	 \loigiai{
		Vì $A+B+C=180^{\circ}$ nên $B+C=180^{\circ}-A$  và $\dfrac{A+B}{2}=\dfrac{180^{\circ}-C}{2}$.\\
		Do  đó
		 \allowdisplaybreaks
		\begin{eqnarray*}
			 &&\cos A+\cos (B+C)+\tan \dfrac{A+B}{2}\\
			&=&\cos A+\cos  \left(180^{\circ}-A\right)+\tan \dfrac{180^{\circ}-C}{2}\\
			&=&\cos A-\cos A+\tan  \left(90^{\circ}-\dfrac{C}{2}\right)\\
			&=&\cot  \dfrac{C}{2}.
		\end{eqnarray*}
	 }
\end{ex}


%%%%%%%%%%%%%%%% TỰ LUẬN %%%%%%%%%%%%%%%%

% nguồn: Ok_TuanLuan_Toan11_Lan2
\begin{ex}%[1D1V3-6]%[ma:327a3d14]
	Chứng minh rằng với mọi tam giác $ABC$, ta đều có  $$\sin A+\sin B+\sin C=4\cos\dfrac{A}{2}\cos\dfrac{B}{2}\cos\dfrac{C}{2}.$$
	\loigiai{ 
	Ta có $ VT=2\sin\dfrac{A+B}{2}\cos\dfrac{A-B}{2}+2\sin\dfrac{C}{2}\cos\dfrac{C}{2}$.\\ 
	Mặt khác, trong tam giác  $ ABC$ ta có$ A+B+C=\pi $ nên $\dfrac{A+B}{2}=\dfrac{\pi}{2}-\dfrac{C}{2}$.\\
	Từ  đó suy ra $\sin\dfrac{A+B}{2}=\cos\dfrac{C}{2}$,$\sin\dfrac{C}{2}=\cos\dfrac{A+B}{2}$.\\
	Vậy 
	\begin{align*}
		VT=&\,2\cos\dfrac{C}{2}\cos\dfrac{A-B}{2}+2\cos\dfrac{A+B}{2}\cos\dfrac{C}{2}\\ 
		 =&\,2\cos\dfrac{C}{2}\left(\cos\dfrac{A-B}{2}+\cos\dfrac{A+B}{2}\right)\\
		 =&\,4\cos\dfrac{C}{2}\cos\dfrac{A}{2}\cos\dfrac{B}{2}\\
		 =&\,VP.
	\end{align*}
	} 
\end{ex}
